\documentclass[../main.tex]{subfiles}

\begin{document}

\chapter{测度论}\label{chapter:appE-测度论}

本附录提供测度论一些基础的快速概述。

\section{Lebesgue 测度}\label{section:appE-Lebesgue测度}

Lebesgue 测度提供了一种描述 $\mathbb{R}^n$ 某些子集的"大小"或"体积"的方法。

\begin{definition}\label{definition:appE-sigma代数}
$\mathbb{R}^n$ 子集的集合 $\mathcal{M}$ 称为 $\sigma$-代数，若
\begin{enumerate}[(i)]
\item $\emptyset,\mathbb{R}^n\in\mathcal{M}$，
\item $A\in\mathcal{M}$ 蕴含 $\mathbb{R}^n\setminus A\in\mathcal{M}$，
\item 若 $\{A_k\}_{k=1}^\infty\subset\mathcal{M}$，则 $\bigcup_{k=1}^\infty A_k,\bigcap_{k=1}^\infty A_k\in\mathcal{M}$。
\end{enumerate}
\end{definition}

\begin{theorem}[Lebesgue 测度和 Lebesgue 可测集的存在性。]\label{theorem:appE-Lebesgue存在}存在 $\mathbb{R}^n$ 子集的 $\sigma$-代数 $\mathcal{M}$ 和映射
\begin{align*}
|\cdot|:\mathcal{M}\to[0,+\infty]
\end{align*}
具有以下性质：
\begin{enumerate}[(i)]
\item $\mathbb{R}^n$ 的每个开子集从而每个闭子集属于 $\mathcal{M}$。
\item 若 $B$ 是 $\mathbb{R}^n$ 中球，则 $|B|$ 等于 $B$ 的 $n$ 维体积。
\item 若 $\{A_k\}_{k=1}^\infty\subset\mathcal{M}$ 且集合 $\{A_k\}_{k=1}^\infty$ 两两不交，则
\begin{align}
\left|\bigcup_{k=1}^\infty A_k\right|=\sum_{k=1}^\infty|A_k| \quad \text{（可数可加性）}.
\label{align:appE-可数可加}
\end{align}
\item 若 $A\subseteq B$，其中 $B\in\mathcal{M}$ 且 $|B|=0$，则 $A\in\mathcal{M}$ 且 $|A|=0$。
\end{enumerate}
\end{theorem}

$\mathcal{M}$ 中的集合称为 Lebesgue 可测集，$|\cdot|$ 是 $n$ 维 Lebesgue 测度。

\begin{remark}\label{remark:appE-测度注}
\begin{enumerate}[(i)]
\item 由 (ii) 和 (iii)，看到 $|A|$ 等于任何具有分段光滑边界的集合 $A$ 的体积。
\item 由\eqref{align:appE-可数可加} 推出
\begin{align}
|\emptyset|=0
\label{align:appE-空集}
\end{align}
和
\begin{align}
\left|\bigcup_{k=1}^\infty A_k\right|\leqslant\sum_{k=1}^\infty|A_k| \quad \text{（可数次可加性）}
\label{align:appE-次可加}
\end{align}
对任何可数可测集集合 $\{A_k\}_{k=1}^\infty$。
\end{enumerate}
\end{remark}

\begin{definition}\label{definition:appE-a.e.}
若某性质在 $\mathbb{R}^n$ 上处处成立，除了一个 Lebesgue 测度为零的可测集外，我们称该性质几乎处处成立，缩写为"a.e."。
\end{definition}

\section{可测函数与积分}\label{section:appE-可测函数}

\begin{definition}\label{definition:appE-可测函数}
设 $f:\mathbb{R}^n\to\mathbb{R}$。称 $f$ 是可测函数，若
\begin{align*}
f^{-1}(U)\in\mathcal{M}
\end{align*}
对每个开子集 $U\subset\mathbb{R}$。
\end{definition}

特别地注意若 $f$ 连续，则 $f$ 可测。两个可测函数的和与积可测。此外若 $\{f_k\}_{k=0}^\infty$ 是可测函数，则 $\limsup f_k$ 和 $\liminf f_k$ 也可测。

\begin{theorem}[Egoroff 定理。]\label{theorem:appE-Egoroff}设 $\{f_k\}_{k=1}^\infty$，$f$ 是可测函数，假设
\begin{align*}
f_k\to f \quad \text{a.e. 在 } A \text{ 上},
\end{align*}
其中 $A\subset\mathbb{R}^n$ 可测，$|A|<\infty$。则对每个 $\varepsilon>0$ 存在可测子集 $E\subset A$ 使得
\begin{enumerate}[(i)]
\item $|A\setminus E|\leqslant\varepsilon$，
\item $f_k\to f$ 在 $E$ 上一致。
\end{enumerate}
\end{theorem}

若 $f$ 是非负可测函数，可用简单函数逼近 $f$ 来定义 Lebesgue 积分
\begin{align*}
\int_{\mathbb{R}^n}f\,\mathrm{d}x.
\end{align*}
见下 \S E.5。若 $f$ 连续或 Riemann 可积，这与通常积分一致。若 $f$ 可测但不一定非负，定义
\begin{align*}
\int_{\mathbb{R}^n}f\,\mathrm{d}x=\int_{\mathbb{R}^n}f^+\,\mathrm{d}x-\int_{\mathbb{R}^n}f^-\,\mathrm{d}x,
\end{align*}
只要右端至少一项有限。此时称 $f$ 可积。

\begin{definition}\label{definition:appE-可和}
可测函数 $f$ 称为可和的，若
\begin{align*}
\int_{\mathbb{R}^n}|f|\,\mathrm{d}x<\infty.
\end{align*}
\end{definition}

仔细注意我们的术语：可测函数若有积分（可等于 $+\infty$ 或 $-\infty$）则可积，若此积分有限则可和。

\begin{definition}\label{definition:appE-本质上确界}
若实值函数 $f$ 可测，定义 $f$ 的本质上确界为
\begin{align*}
\operatorname*{ess\,sup}f\coloneqq\inf\{\mu\in\mathbb{R}\mid|\{f>\mu\}|=0\}.
\end{align*}
\end{definition}

\section{积分收敛定理}\label{section:appE-积分收敛}

Lebesgue 积分论特别有用，因它提供以下强大的收敛定理。

\begin{theorem}[Fatou 引理。]\label{theorem:appE-Fatou}假设函数 $\{f_k\}_{k=1}^\infty$ 非负可测。则
\begin{align*}
\int_{\mathbb{R}^n}\liminf_{k\to\infty}f_k\,\mathrm{d}x\leqslant\liminf_{k\to\infty}\int_{\mathbb{R}^n}f_k\,\mathrm{d}x.
\end{align*}
\end{theorem}

\begin{theorem}[单调收敛定理。]\label{theorem:appE-单调收敛}假设函数 $\{f_k\}_{k=1}^\infty$ 可测，
\begin{align*}
0\leqslant f_1\leqslant f_2\leq\cdots\leqslant f_k\leqslant f_{k+1}\leqslant\cdots.
\end{align*}
则
\begin{align*}
\int_{\mathbb{R}^n}\lim_{k\to\infty}f_k\,\mathrm{d}x=\lim_{k\to\infty}\int_{\mathbb{R}^n}f_k\,\mathrm{d}x.
\end{align*}
\end{theorem}

\begin{theorem}[控制收敛定理。]\label{theorem:appE-控制收敛}假设函数 $\{f_k\}_{k=1}^\infty$ 可积且
\begin{align*}
f_k\to f \quad \text{a.e.}
\end{align*}
又假设
\begin{align*}
|f_k|\leqslant g \quad \text{a.e.}
\end{align*}
对某可和函数 $g$。则
\begin{align*}
\int_{\mathbb{R}^n}f_k\,\mathrm{d}x\to\int_{\mathbb{R}^n}f\,\mathrm{d}x.
\end{align*}
\end{theorem}

\section{微分}\label{section:appE-微分}

一个重要事实是可和函数在几乎每点"近似连续"。

\begin{theorem}[Lebesgue微分定理]\label{theorem:appE-Lebesgue微分}
设 $f:\mathbb{R}^n\to\mathbb{R}$ 局部可和。
\begin{enumerate}[(i)]
\item 对几乎处处的点 $x_0\in\mathbb{R}^n$，当$r\to 0$时,有
\begin{align}
\fint_{B(x_0,r)}f\,\mathrm{d}x\to f(x_0) \quad \text{当 } r\to 0
\label{align:appE-平均收敛}
\end{align}

\item 事实上，对几乎处处的点 $x_0\in\mathbb{R}^n$，当$r\to 0$时,有
\begin{align}
\fint_{B(x_0,r)}|f(x)-f(x_0)|\,\mathrm{d}x\to 0 .
\label{align:appE-Lebesgue点}
\end{align}
\end{enumerate}
称满足\eqref{align:appE-Lebesgue点}式的点 $x_0$ 称为 $f$ 的 Lebesgue 点。
\end{theorem}

\begin{theorem}[Lebesgue微分定理]\label{theorem:appE-Lebesgue微分1}
若对某 $1\leqslant p<\infty$，$f\in L^p_{\mathrm{loc}}(\mathbb{R}^n)$,则对几乎处处的点 $x_0\in\mathbb{R}^n$ ，当$r\to 0$时,有
\begin{align*}
\fint_{B(x_0,r)}|f(x)-f(x_0)|^p\,\mathrm{d}x\to 0.
\end{align*}
\end{theorem}

\section{Banach 空间值函数}\label{section:appE-Banach值函数}

我们将可测性、可积性等概念推广到映射
\begin{align*}
\mathbf{f}:[0,T]\to X
\end{align*}
其中 $T>0$，$X$ 是实 Banach 空间，范数为 $\|\cdot\|$。

\begin{definition}\label{definition:appE-强可测}
\begin{enumerate}[(i)]
\item 函数 $\mathbf{s}:[0,T]\to X$ 称为简单的，若它具有形式
\begin{align}
\mathbf{s}(t)=\sum_{i=1}^m\chi_{E_i}(t)u_i \quad (0\leqslant t\leqslant T),
\label{align:appE-简单函数}
\end{align}
其中每个 $E_i$ 是 $[0,T]$ 的 Lebesgue 可测子集，$u_i\in X$（$i=1,\cdots,m$）。

\item 函数 $\mathbf{f}:[0,T]\to X$ 强可测，若存在简单函数 $\mathbf{s}_k:[0,T]\to X$ 使得
\begin{align*}
\mathbf{s}_k(t)\to\mathbf{f}(t)
\end{align*}
对 a.e. $0\leqslant t\leqslant T$。

\item 函数 $\mathbf{f}:[0,T]\to X$ 弱可测，若对每个 $u^*\in X^*$，映射 $t\mapsto\langle u^*,\mathbf{f}(t)\rangle$ 是 Lebesgue 可测的。
\end{enumerate}
\end{definition}

\begin{definition}\label{definition:appE-几乎可分}
称 $\mathbf{f}:[0,T]\to X$ 几乎可分取值，若存在子集 $N\subset[0,T]$，$|N|=0$，使得集合 $\{\mathbf{f}(t)\mid t\in[0,T]\setminus N\}$ 可分。
\end{definition}

\begin{theorem}[强可测与弱可测。]\label{theorem:appE-强弱可测}映射 $\mathbf{f}:[0,T]\to X$ 强可测当且仅当 $\mathbf{f}$ 弱可测且几乎可分取值。
\end{theorem}

\begin{definition}\label{definition:appE-Bochner积分}
\begin{enumerate}[(i)]
\item 若 $\mathbf{s}(t)=\sum_{i=1}^m\chi_{E_i}(t)u_i$ 简单，定义
\begin{align}
\int_0^T\mathbf{s}(t)\,\mathrm{d}t\coloneqq\sum_{i=1}^m|E_i|u_i.
\label{align:appE-简单积分}
\end{align}

\item 称强可测函数 $\mathbf{f}:[0,T]\to X$ 可和，若存在简单函数序列 $\{\mathbf{s}_k\}_{k=1}^\infty$ 使得
\begin{align}
\int_0^T\|\mathbf{s}_k(t)-\mathbf{f}(t)\|\,\mathrm{d}t\to 0 \quad \text{当 } k\to\infty.
\label{align:appE-可和条件}
\end{align}

\item 若 $\mathbf{f}$ 可和，定义
\begin{align}
\int_0^T\mathbf{f}(t)\,\mathrm{d}t=\lim_{k\to\infty}\int_0^T\mathbf{s}_k(t)\,\mathrm{d}t.
\label{align:appE-Bochner积分}
\end{align}
\end{enumerate}
\end{definition}

\begin{theorem}[可和函数。]\label{theorem:appE-可和函数}强可测函数 $\mathbf{f}:[0,T]\to X$ 可和当且仅当 $t\mapsto\|\mathbf{f}(t)\|$ 可和。此时
\begin{align*}
\left\|\int_0^T\mathbf{f}(t)\,\mathrm{d}t\right\|\leqslant\int_0^T\|\mathbf{f}(t)\|\,\mathrm{d}t,
\end{align*}
且
\begin{align*}
\left\langle u^*,\int_0^T\mathbf{f}(t)\,\mathrm{d}t\right\rangle=\int_0^T\langle u^*,\mathbf{f}(t)\rangle\,\mathrm{d}t
\end{align*}
对每个 $u^*\in X^*$。
\end{theorem}

测度论的好书有 Folland~\cite{F2} 和 DiBenedetto~\cite{DB2}。\cref{theorem:appE-强弱可测}、\cref{theorem:appE-可和函数} 的证明见 Yosida~\cite{Y} 第 V 章 \S4--5。

\end{document}